% Generated by make_paper_v2.py from paper_v2_template.tex + results_v2 numbers.
% Do not edit numbers by hand: edit the template and refill.
\documentclass{ieeeaccess}

\IEEEoverridecommandlockouts
\doi{}
\IEEEpubid{\begin{minipage}{\textwidth}\footnotesize This work is licensed under a Creative Commons Attribution 4.0 License. For more information, see \href{https://creativecommons.org/licenses/by/4.0/}{https://creativecommons.org/licenses/by/4.0/}.\end{minipage}}

\usepackage{amsmath}
\usepackage{booktabs}
\usepackage{graphicx}
\usepackage[colorlinks=true]{hyperref}

\begin{document}

\graphicspath{{./figs/}{../figs_journal_clean/}}
\history{Date of publication xxxx 00, 0000, date of current version xxxx 00, 0000.}
\doi{10.1109/ACCESS.2026.XXXXXXX}
\title{Multi-Horizon Failure-Risk Models for Battery Failure Prediction: How Much Cross-Chemistry Transfer Is a State-of-Health Shortcut?}
\author{\uppercase{Anonymous}\authorrefmark{1}, \uppercase{Anonymous}\authorrefmark{2}, \uppercase{Anonymous}\authorrefmark{3}, AND \uppercase{Anonymous}\authorrefmark{4}}
\address[1]{Affiliation 1, City, Country (e-mail: author1@example.com)}
\address[2]{Affiliation 2, City, Country (e-mail: author2@example.com)}
\address[3]{Affiliation 3, City, Country (e-mail: author3@example.com)}
\address[4]{Affiliation 4, City, Country (e-mail: author4@example.com)}
\tfootnote{This work received no external funding. The authors declare no conflicts of interest.}
\markboth{Anonymous \headeretal: Multi-Horizon Failure-Risk Models for Battery Failure Prediction}{Anonymous \headeretal: Multi-Horizon Failure-Risk Models for Battery Failure Prediction}
\corresp{Corresponding author: Anonymous (e-mail: corresponding.author@example.com).}
\emergencystretch=6em
\begin{abstract}
Predicting whether a lithium-ion cell will fail inside a short operating window is a different question from estimating its remaining useful life, and it is the one that matters most for real-time dispatch in electric vehicles and grid storage. This study widens a recently proposed multi-horizon failure-risk classification framework, first shown on one model and one dataset, into a validity-focused evaluation covering four model families, a true discrete-time hazard model, trivial health-state baselines, and @@N_TARGETS@@ lithium-iron-phosphate (LFP) transfer targets, all under cell-disjoint splits and a leakage-clean, cross-fitted calibration protocol. Within-dataset discrimination is reliable, with fold-mean area under the curve (AUC) of @@WITHIN_MIN@@--@@WITHIN_MAX@@ across the tree ensembles under Platt calibration, and Platt better preserves discrimination than isotonic while achieving comparable calibration error. The central result is a controlled State-of-Health (SOH) ablation: apparent cross-dataset, cross-chemistry transfer with SOH included (pooled AUC up to @@XFER_WITH_MAX@@) collapses without SOH (@@XFER_NO_MAX@@ at best), and much of the with-SOH signal is reproducible by a distance-to-threshold rule on SOH alone (@@BASE_DIST_ALL_SEV@@ AUC on the @@N_SEV@@-cell MIT--Stanford Severson set), while the remaining sensor features can even invert under shift. Because SOH also enters the failure definition, all with-SOH numbers are upper bounds that partly measure label proximity rather than degradation knowledge: relabeling failures by voltage sag alone, which SOH does not define, removes most of the with-SOH advantage. Same-chemistry cross-laboratory controls confirm that the collapse is driven by dataset shift at least as much as by chemistry, and SOH's dominance is therefore a dataset- and chemistry-dependent health-state shortcut rather than a chemistry-specific transferable signal. The truncation-prone isotonic calibrator largely fails to transfer, cell-level bootstrap intervals replace cycle-level significance tests as the primary evidence, and the complete deployment side is retained: three tree ensembles (900 trees, 26~336 nodes) fit into a 372~kB binary and run on a \$12 ESP32-S3 in 0.32--2.36~ms per inference, reproducing the Python training libraries to within $1.1\times10^{-6}$ across all 1~028 validation rows.
\end{abstract}
\begin{keywords}
Battery failure prediction, multi-horizon failure-risk classification, state-of-health shortcut, cross-chemistry transfer, dataset shift, probability calibration, discrete-time hazard model, SHAP, embedded machine learning, ESP32-S3.
\end{keywords}
\titlepgskip=-21pt
\maketitle
\emergencystretch=6em
\makeatletter\def\UrlBreaks{\do\/\do\-}\makeatother
\section{Introduction}
Lithium-ion batteries sit underneath the electrification of transport and the decarbonization of grids, yet their long-term reliability remains hard to certify in the field. Conventional prognostics casts the problem as remaining-useful-life (RUL) regression, in which a model predicts a continuous cycle count to end of life \cite{meng2019review}. RUL regression is appropriate for capacity-fade studies, but it answers a question operators rarely ask in real time: a battery management system needs to know whether a specific cell will survive the next mission, not how many cycles it has left in aggregate. Review studies group prior methods into physics-based, purely data-driven, and hybrid families \cite{meng2019review}. Tree ensembles, such as Random Forest \cite{breiman2001random}, XGBoost \cite{chen2016xgboost}, and LightGBM \cite{ke2017lightgbm}, have been applied to capacity and RUL estimation, but their use for binary failure-risk classification at multiple horizons is recent \cite{shikdar2026learning}.

Shikdar and Laaksonen \cite{shikdar2026learning} recast the problem as multi-horizon failure-risk classification. For a given cycle $t$ and horizon $H$, the model predicts whether failure, defined as SOH below 0.80 or average voltage sag beyond 6\%, occurs within the next $H$ cycles. Their histogram-based gradient boosting model on the NASA 18650 set (37 LCO cells) reached AUC of 0.87--0.90 across $H \in \{10,20,30,50\}$. That work left three gaps open. First, one model on one dataset says nothing about sensitivity to the learner, the hyperparameters, or the calibration method. Second, no cross-dataset check existed, so nothing stopped a model trained on LCO cells from being trusted on LFP cells. Third, no path existed from a Python research prototype to the microcontroller where a battery management system actually runs.

This paper begins from those gaps but is organized around a sharper validity question that emerged while closing them: \emph{when a failure-risk model appears to transfer between datasets and chemistries, is anything about degradation being transferred, or is the model merely reading a health-state variable that encodes proximity to the failure criterion itself?} The question matters because the composite failure label is defined from SOH, and SOH is also the strongest available feature. A model can therefore score well on transfer while learning ``distance to the SOH $=$ 0.80 threshold,'' a quantity that says nothing about chemistry-specific degradation physics. A hostile reading of any cross-chemistry result that ignores this confound is, in our experience, correct more often than not.

Two adjacent literatures bear on these questions. Transfer learning for battery state estimation is a growing area: Sahoo et al. \cite{sahoo2022transfer} fine-tuned an SOH estimator across NASA and CALCE cells, and Lu et al. \cite{lu2023deep} pretrained on large-scale degradation data before adapting to new cell types. Both target regression-based SOH estimation rather than classification-based failure risk, and the classification setting adds two complications: the label depends on the chosen horizon $H$, and the calibration of probability outputs, which decisions depend on, may degrade under distribution shift. On calibration, Platt scaling \cite{platt1999probabilistic} fits a sigmoid through logistic regression on the model's raw scores, while isotonic regression \cite{zadrozny2002transforming} fits a non-decreasing step function. Niculescu-Mizil and Caruana \cite{niculescu2005predicting} showed that Platt is safer on small datasets and isotonic can overfit, and Huang et al. \cite{huang2020experimental} extended this to imbalanced classification; no prior study has compared the two calibrators under cross-chemistry distribution shift in battery prognostics. A companion robustness study further found that under realistic operating perturbation the calibration layer is the primary point of failure, and that recalibration on a small sample of the operational data largely restores calibration \cite{shikdar2026robustness}; whether the same repair survives a cross-chemistry shift is examined in this work. On explainability and deployment, TreeSHAP \cite{lundberg2020from} provides consistent feature attributions and has become the standard tool for diagnosing ensemble inputs, while edge inference removes latency, connectivity requirements, and privacy concerns \cite{grzesik2024combining}. Online recalibration methods for shifting score distributions are also being studied \cite{gupta2023online}. Table~\ref{tab:lit} situates this work against these lines of research.

\begin{table*}[!t]
\caption{Comparison of Representative Battery Studies With This Work}
\label{tab:lit}
\centering
\resizebox{\textwidth}{!}{%
\begin{tabular}{lllp{0.27\textwidth}}
\toprule
\textbf{Study} & \textbf{Family} & \textbf{Data} & \textbf{Position relative to this work}\\
\midrule
Meng and Li \cite{meng2019review} & Prognostics review & Various & Groups methods into physics-based, data-driven, and hybrid families\\
Breiman \cite{breiman2001random}, Chen and Guestrin \cite{chen2016xgboost}, Ke et al. \cite{ke2017lightgbm} & Tree ensembles & Capacity, RUL & Established ensembles for RUL regression, not failure-risk classification\\
Sahoo et al. \cite{sahoo2022transfer} & Transfer learning & NASA, CALCE & Transferred SOH regression across chemistries; regression only\\
Lu et al. \cite{lu2023deep} & Deep learning & Large SOH corpora & Pretraining for SOH estimation; regression only\\
Platt \cite{platt1999probabilistic}, Zadrozny and Elkan \cite{zadrozny2002transforming}, Niculescu-Mizil and Caruana \cite{niculescu2005predicting}, Huang et al. \cite{huang2020experimental} & Probability calibration & Various & Calibrators compared in-domain; no cross-dataset setting\\
Lundberg et al. \cite{lundberg2020from} & Explainability & Various & TreeSHAP attributions for tree ensembles\\
Grzesik and Mrozek \cite{grzesik2024combining} & Edge computing & Various & Edge inference review without battery failure use cases\\
Gupta and Ramdas \cite{gupta2023online} & Calibration & Various & Online Platt scaling under distribution shift\\
Shikdar and Laaksonen \cite{shikdar2026learning} & Failure-risk classification & NASA & Single model, single dataset; baseline of this paper\\
\textbf{This work} & Failure-risk classification + validity analysis & NASA, CALCE, Oxford, Severson & Multi-model benchmark, SOH-shortcut diagnosis with label-, feature-, and dataset-level controls, discrete-time hazard model, calibration under shift, embedded deployment\\
\bottomrule
\end{tabular}}
\end{table*}

The objective of this study is to close the three gaps identified above and, in doing so, to answer the validity question rigorously. The contributions are:
\begin{itemize}
\item a four-model benchmark, three tree ensembles (XGBoost, LightGBM, and Random Forest) plus a GRU sequence classifier, run on two LCO datasets (NASA 18650 and CALCE LCO/CX2) under cell-disjoint GroupKFold and leave-one-cell-out protocols with cross-fitted, leakage-free calibration;
\item a true discrete-time hazard formulation, in which per-offset hazards are pooled into a survival product that is monotone in the horizon by construction, evaluated against the four independent fixed-horizon classifiers, whose predicted risks violate the corresponding monotonicity requirement in @@MONO_RATE_WITHIN@@ of within-dataset cases;
\item a shortcut diagnosis of cross-dataset, cross-chemistry transfer: minimal baselines (a distance-to-threshold rule, SOH-only, cycle-only, and sensor-only logistic regressions), a factorial feature ablation over \{SOH, cycle index, sensors\}, a failure-definition ablation that relabels failures by SOH or by voltage sag alone, and same-chemistry cross-laboratory controls (NASA$\leftrightarrow$CALCE and Severson$\to$Oxford) that separate dataset shift from chemistry shift;
\item an uncertainty treatment in which the cell, not the cycle, is the experimental unit: cell-level bootstrap confidence intervals on every headline number, with cycle-level DeLong tests demoted to secondary evidence and their pseudoreplication caveat stated explicitly;
\item a reliability analysis of probability calibration under shift --- Platt, isotonic, and temperature scaling, all fitted on out-of-fold source scores --- including recalibration on small labeled LFP samples, and a cost-sensitive operational evaluation with source-selected decision thresholds;
\item a complete embedded deployment that packs all three ensembles into one flat binary and runs them on an ESP32-S3 through a hand-written C tree walker.
\end{itemize}

The remainder of this paper is organized as follows. Section~II describes the methodology, including the failure labels and their endpoints, the datasets, the models including the discrete-time hazard formulation, the calibration and transfer protocols, and the deployment architecture. Section~III presents the experimental results. Section~IV discusses the analysis, implications, and limitations, and Section~V concludes.

\section{Methodology}

\subsection{Composite Failure Label}
Following the original protocol \cite{shikdar2026learning}, the State-of-Health of cell $c$ at cycle $k$ is the ratio of the discharge capacity to the mean capacity over the first 10 cycles,
\begin{equation}\label{eq:soh}
\begin{aligned}
\mathrm{SOH}_k^{(c)} &= \frac{Q_k^{(c)}}{\bar{Q}^{(c)}},\\[2pt]
\bar{Q}^{(c)} &= \frac{1}{10}\sum_{i=1}^{10} Q_i^{(c)},
\end{aligned}
\end{equation}
where $Q_k^{(c)}$ is the discharge capacity at cycle $k$. The early-life voltage baseline is defined analogously, $\bar{V}^{(c)}=\frac{1}{10}\sum_{i=1}^{10}\bar{V}_i^{(c)}$, with $\bar{V}_i^{(c)}$ the average discharge voltage. A discharge cycle $t$ in cell $c$ is labeled as failing within horizon $H$ if any cycle $k\in[t,t{+}H)$ meets either of two criteria,
\begin{equation}\label{eq:label}
y_t^{(c,H)} \;=\; \begin{cases}
1, & \exists\, k\in[t,t{+}H):\ \mathrm{SOH}_k^{(c)}\le 0.80 \ \text{or}\\[2pt]
& \bar{V}_k^{(c)} < 0.94\,\bar{V}^{(c)},\\[2pt]
0, & \text{otherwise,}
\end{cases}
\end{equation}
and once triggered the label stays 1 for all later cycles. The task is to estimate the cumulative failure risk over the horizon,
\begin{equation}\label{eq:risk}
p_t^{(c,H)} \;=\; \Pr\!\bigl(T_f \le t{+}H \mid T_f > t,\ \mathbf{x}_t^{(c)}\bigr),
\end{equation}
where $T_f$ is the failure cycle and $\mathbf{x}_t^{(c)}$ is the feature vector measured at cycle $t$. Datasets without voltage data, or where voltage sits at the cutoff as in Oxford's flat LFP plateau, fall back to the SOH criterion alone.

Three properties of this label definition matter for everything that follows. First, (3) is a fixed-horizon cumulative failure risk, not a discrete-time hazard: each horizon is an independent binary classifier, no cycle-level hazard $h_k$ or survival product is estimated, and no monotonicity constraint $p_t^{(c,10)}\le p_t^{(c,20)}\le p_t^{(c,30)}\le p_t^{(c,50)}$ is enforced. Section~II-G adds a true discrete-time hazard model, and Section~III-H quantifies how often the four independent classifiers violate monotonicity. Second, the label is \emph{partially defined by SOH itself}, and SOH is also an input feature; a with-SOH model can therefore detect proximity to the SOH $=0.80$ label threshold without learning any degradation dynamics. We treat every with-SOH result in this paper as an upper bound that includes this label-proximity leakage, and Section~III-F measures it directly by relabeling failures with a single endpoint at a time. Third, the two endpoints can be separated: the label can be defined by the SOH criterion alone, by the voltage-sag criterion alone, or by both, and the model's feature set can be ablated independently of the label definition.

\subsection{Datasets}
Four public aging datasets are used:
\begin{itemize}
\item \textbf{1) NASA 18650} \cite{saha2007battery}: 37 LCO cells aged under randomized charge--discharge profiles at room temperature, roughly 1~000 cycles per cell, with diverse failure modes;
\item \textbf{2) CALCE LCO/CX2} \cite{calce2023battery}: 7 cells (CS2\_33--36, CX2\_36--38) at 1C/1C with per-cycle voltage, current, and capacity; per-cell ambient temperature and discharge duration are unlogged;
\item \textbf{3) Oxford} \cite{birkl2017oxford}: 5 LFP pouch cells at 1C/1C under controlled temperature. The public release stores discharge summaries at roughly 100-cycle intervals, so after cleaning, 46--78 usable discharge cycles per cell remain (320 rows total);
\item \textbf{4) MIT--Stanford Severson} \cite{severson2019data}: 141 LFP cells under a fast-charging protocol, 534--2~237 cycles each (about 117~000 in total).
\end{itemize}

One consequence of the Oxford logging resolution deserves emphasis. Because consecutive cleaned cycles are roughly 100 cycles apart, every horizon $H\in\{10,20,30,50\}$ is shorter than the sampling interval, and the set of positive-label rows on Oxford is \emph{identical for all four horizons}; Oxford scores therefore measure pooled failure detection rather than horizon-resolved prediction, and AUC differences between horizons on Oxford come entirely from the source-side model, not from the target labels. Per-cell statistics on Oxford effectively cover four of the five cells, because one cell contributes too few positive rows for a within-cell AUC to be defined. Both caveats apply to every Oxford number in this paper.

\subsection{Features and Preprocessing}
Each cycle is described by seven scalar features: cycle index, average discharge voltage, minimum discharge voltage, average discharge current, average temperature, discharge duration, and SOH. Trees are scale-invariant, so no standardization is applied to them; the GRU and all logistic-regression baselines receive per-feature standardization using training-set statistics. Missing CALCE values (temperature and duration are unlogged) are filled with zero. Within a single dataset a constant feature cannot be used for splitting, but in pooled cross-dataset training a constant zero column acts as a dataset identifier; Section~III-G quantifies this confound with a common-feature control.

\subsection{Tree Models and Hyperparameters}
Three tree ensembles are benchmarked with hyperparameters matched to the original study \cite{shikdar2026learning}: XGBoost \cite{chen2016xgboost} with max\_depth 4, 300 estimators, learning rate 0.05, subsample 0.8, colsample\_bytree 0.8, and min\_child\_weight 5; LightGBM \cite{ke2017lightgbm} with the same depth, count, and rate; and Random Forest \cite{breiman2001random} with max\_depth 6 and 300 estimators. All use random\_state 42 for deterministic replication.

\subsection{GRU Sequence Classifier}
A compact GRU tests whether sequence-aware representations capture trajectories that transfer across chemistries. It ingests a sliding window of $W{=}10$ consecutive cycles per cell as a $10{\times}f$ tensor and runs one GRU layer with 8 hidden units,
\begin{equation}\label{eq:gru}
\begin{aligned}
\mathbf{r}_t &= \sigma(W_r\mathbf{x}_t + U_r\mathbf{h}_{t-1}),\\[2pt]
\mathbf{z}_t &= \sigma(W_z\mathbf{x}_t + U_z\mathbf{h}_{t-1}),\\[2pt]
\tilde{\mathbf{h}}_t &= \tanh(W_h\mathbf{x}_t + U_h(\mathbf{r}_t\odot\mathbf{h}_{t-1})),\\[2pt]
\mathbf{h}_t &= (1-\mathbf{z}_t)\odot\mathbf{h}_{t-1} + \mathbf{z}_t\odot\tilde{\mathbf{h}}_t,
\end{aligned}
\end{equation}
and the final hidden state maps through a linear layer with sigmoid activation, $p=\sigma(\mathbf{w}^{\top}\mathbf{h}_W + b)$, to give the failure probability. Adam with learning rate 0.005 and binary cross-entropy loss are used. The 8-unit width was chosen to match the small number of training cells (7 to 37); wider layers overfit severely. For cross-dataset runs the feature set is reduced to the common, unit-safe set --- cycle number, average voltage, minimum voltage, and SOH when enabled --- standardized with training-set statistics; average current is dropped (amperes for LCO sources, derivative-based estimates for the LFP targets), as are temperature and duration. The trees are additionally evaluated on this identical common feature set (Section~III-C), so that architecture and feature availability are not confounded: where the GRU and the trees use different inputs, we say so explicitly, and no architecture claim rests on a feature-confounded comparison.

\subsection{Minimal Baselines}
Because the central question is whether apparent performance is a health-state shortcut, the most important models are embarrassingly simple ones \cite{shikdar2026learning}:
\begin{itemize}
\item \textbf{SOH distance rule}: the fitted-free score $d_t = 0.80 - \mathrm{SOH}_t$, i.e., distance below the label threshold;
\item \textbf{SOH-only} logistic regression, $\Pr(y_t^{(H)}{=}1)=\sigma(a_H\,\mathrm{SOH}_t+b_H)$;
\item \textbf{cycle-only} logistic regression on the cycle index;
\item \textbf{SOH$+$cycle} logistic regression;
\item \textbf{sensors-only} logistic regression on the five non-SOH, non-cycle features;
\item \textbf{full-feature} logistic regression on all seven features.
\end{itemize}
If the SOH-only and distance-rule baselines reproduce most of the transferable performance of the tuned ensembles, then the ensembles contribute little beyond reading the health state --- which is precisely the negative result the validity question predicts, and the one an evaluation-oriented paper should surface.

\subsection{A True Discrete-Time Hazard Model}
To complement the four independent fixed-horizon classifiers, we also fit an explicit discrete-time survival model. For each observed cycle $t$ and each offset $\ell\in\{1,\dots,L\}$ with $L{=}50$, define the hazard
\begin{equation}\label{eq:hazard}
h_\ell(t) \;=\; \Pr\!\bigl(T_f = t{+}\ell \mid T_f \ge t{+}\ell,\ \mathbf{x}_t\bigr),
\end{equation}
conditional on the features measured at $t$ only, so the formulation remains operational (no future measurements are required). One pooled classifier --- XGBoost with the same hyperparameters as above, and a logistic regression --- is fit on the offset-expanded at-risk training rows $(\mathbf{x}_t, \ell)$ with target $\mathbf{1}\{T_f = t{+}\ell\}$, restricted to rows with $t{+}\ell \le T_f$, and with $\ell$ as an additional input feature. The cumulative risk over horizon $H$ is the survival product
\begin{equation}\label{eq:survival}
R(t,H) \;=\; 1 - \prod_{\ell=1}^{H}\bigl(1-\hat h_\ell(t)\bigr),
\end{equation}
which is monotone non-decreasing in $H$ by construction and yields all four horizons from one model. Rows at or after failure are scored exactly like the fixed-horizon classifiers score them, so the evaluation rows are identical across formulations.

\subsection{Calibration Protocol}
Three post-hoc calibration schemes are compared: Platt scaling fits a sigmoid to the base model's raw scores $s_i$,
\begin{equation}\label{eq:platt}
p_i \;=\; \frac{1}{1+\exp(-\alpha s_i - \beta)},
\end{equation}
by logistic regression; isotonic regression \cite{zadrozny2002transforming} fits the non-decreasing step function that minimizes squared error, implemented with scikit-learn's \texttt{IsotonicRegression} and out-of-bounds clipping; and temperature scaling divides the logit of the score by a single fitted temperature.

The protocol is leakage-clean. For every within-dataset evaluation fold, the calibrators are fitted on \emph{cross-fitted out-of-fold scores}: an inner cell-disjoint GroupKFold over the training cells produces, for every training row, a score from a model that never saw that row's cell, and the calibrators are fitted on those out-of-fold scores; the untouched test cells are scored by a model trained on all training cells and calibrated with the out-of-fold-fitted functions. In the transfer setting the transfer model trains on all source cells, the calibrators are fitted on source out-of-fold scores, and the target cells are scored and calibrated without ever entering any fit. No calibrator in this paper is ever fitted on the scores it is applied to, and no test cell contributes to training or calibration. The earlier version of this study fitted calibrators on in-sample training scores; the cross-fitted absolute calibration numbers reported here are accordingly more conservative.

A reproducibility note is retained from the original study comparison. The original study \cite{shikdar2026learning} reported Brier scores around 0.032; re-execution of the released code yields Brier values between roughly 0.09 and 0.29, a gap of up to ninefold that could not be resolved from the available source. The AUC values reported here are comparable to the published range, and the qualitative conclusions hold. The Brier gap does not affect the relative comparisons that form the core of this paper.

\subsection{Transfer Protocols}
Transfer is evaluated by training on LCO cells (NASA alone, CALCE alone, or both) and testing on the cells of the target datasets. Performance is reported per cell (mean $\pm$ standard deviation across test cells) and pooled over all target rows with cell-level bootstrap confidence intervals; pooling alone can hide degenerate behavior where a model assigns nearly identical scores to every cell, and per-cell statistics expose it. Three families of transfer are used:

\begin{itemize}
\item \textbf{Cross-chemistry} (LCO$\to$LFP): NASA, CALCE, or both $\to$ Oxford and $\to$ Severson, with and without SOH as a feature, and on the common feature set shared with the GRU;
\item \textbf{Same-chemistry, cross-laboratory controls}: NASA$\to$CALCE and CALCE$\to$NASA (LCO$\to$LCO) and Severson$\to$Oxford and Oxford$\to$Severson (LFP$\to$LFP). Because chemistry, laboratory, cell manufacturer, cycling protocol, temperature, and logging resolution all change simultaneously between any source and target, the cross-chemistry shifts above are cross-\emph{dataset} shifts as well; ``cross-chemistry'' is used throughout as shorthand, and these controls are what allow dataset shift and chemistry shift to be told apart;
\item \textbf{Feature-definition ablations}: the factorial feature grid of Section~III-E and the failure-endpoint relabeling of Section~III-F.
\end{itemize}

\subsection{Evaluation and Statistical Testing}
\label{sec:eval}
Within-dataset evaluation uses 5-fold cell-grouped cross-validation on NASA (37 cells) and leave-one-cell-out on CALCE (7 cells), so every cycle from a given cell stays entirely on one side of every split --- train, calibration, and test are cell-disjoint by construction, and models are retrained from scratch per horizon. The experimental unit is the cell, not the cycle: cycles within a cell are strongly correlated, so headline AUCs are reported as fold means with fold standard deviations, and every headline number carries a 95\% bootstrap confidence interval obtained by resampling \emph{cells} with replacement (2~000 draws, or 400 for the 117~000-row Severson pools) and recomputing the statistic. Paired comparisons --- above all with-SOH versus without-SOH --- use the same cell resamples for both models, giving a paired bootstrap interval on the AUC difference as the primary evidence.

The DeLong nonparametric test \cite{delong1988comparing} for paired ROC curves is retained as secondary evidence. It evaluates
\begin{equation}\label{eq:delong}
z \;=\; \frac{A_1 - A_2}{\sqrt{V_{11}+V_{22}-2V_{12}}},
\end{equation}
where $A_j$ is the empirical area under the $j$-th ROC curve and $V_{ij}$ are covariance terms estimated from the same test instances. The test treats every cycle-level prediction as an independent observation; because cycles within a cell are correlated, the effective sample size is far smaller than the row count, and the resulting $p$-values are inflated by pseudoreplication. We therefore report them with that caveat and read them only as evidence that an effect is directionally consistent, never as precise significance statements.

Discrimination alone is not enough for a reliability-oriented study, so each within-dataset configuration also reports the Brier score, expected calibration error (10 equal-width bins), log loss, the slope and intercept of a logistic recalibration of the predicted probabilities (slope 1 and intercept 0 indicate perfect calibration), the area under the precision--recall curve, which is sensitive to the horizon-dependent class prevalence that AUC ignores, and sensitivity at a fixed 10\% false-positive rate. Operational, cost-sensitive evaluation is described with the results in Section~III-K.

\subsection{Embedded Deployment Architecture}
The deployment pipeline exports the three trained models into one flat binary and runs them on an ESP32-S3 through a hand-written C tree walker. The layout is deliberately simple: a 12-byte file header, a $3{\times}4$-byte offset table, and per-model sections holding a 4-byte tree count followed by consecutive per-tree node arrays. Each node stores the split feature index, the split threshold as a 32-bit float, and child indices; leaves store the class-1 probability directly. The walker is one C function of about 70 lines that follows child pointers, sums the leaf contributions per tree, and applies the per-model output transform. Writing $\ell_t$ for the leaf value reached in tree $t$ and $T$ for the number of trees, Random Forest averages and clamps its leaf probabilities,
\begin{equation}\label{eq:rf}
p \;=\; \mathrm{clamp}\!\left(\frac{1}{T}\sum_{t=1}^{T}\ell_t,\ 0,\ 1\right),
\end{equation}
while XGBoost and LightGBM pass the summed score through a sigmoid,
\begin{equation}\label{eq:sigmoid}
p \;=\; \frac{1}{1+\exp\!\bigl(-\textstyle\sum_{t=1}^{T}\ell_t - s_0\bigr)},
\end{equation}
where $s_0$ is the model's initialization score ($s_0\neq 0$ only for XGBoost). Fig.~\ref{fig:circuit} shows the wiring. Three validation stages confirm correctness: a Python walker that mirrors the C logic and is compared against \texttt{predict\_proba()}, a cross-compiled C binary tested on desktop x86\_64, and on-target execution on an ESP32-S3-DevKitC-1 reading the binary from flash. All three must report maximum absolute error below $10^{-5}$ on every row before deployment is accepted.

\begin{figure*}[!t]
\centering
\includegraphics[width=0.85\textwidth]{circuit.png}
\caption{Prototype wiring for the ESP32-S3 deployment. The battery discharges through the INA219 shunt into a load, the 2:1 divider senses cell voltage at ADC GPIO1, and I2C (GPIO8/9) serves the INA219 and DS3231 while OneWire (GPIO10) reads the DS18B20. The boards are a \$12 ESP32-S3-DevKitC-1 and two breakout modules.}
\label{fig:circuit}
\end{figure*}
